% --- INICIO DE PLANTILLA PERSONALIZADA ---
\documentclass[oneside,11pt]{Latex/Classes/PhDthesisPSnPDF}

% --- Paquetes Originales ---
\usepackage{blindtext}
\usepackage{actuarialangle} 
\usepackage{tabularx}
\usepackage{float}
\usepackage{multirow, array}
\usepackage[usenames,dvipsnames,svgnames,table]{xcolor}
\usepackage{longtable}
\usepackage{latexsym,amsmath,amssymb,amsfonts}
\usepackage{enumitem}
\usepackage{xparse}
\usepackage{booktabs}
\usepackage[round, sort, numbers]{natbib}  
\usepackage{adjustbox}
\usepackage[utf8]{inputenc}
\usepackage{caption}
\usepackage{graphicx}
\usepackage{hyperref}
\usepackage{cleveref}

% Definición de colores
\definecolor{miblue}{RGB}{68, 114, 196}

% --- CAMBIO CRÍTICO 1: Usar \input para macros en el preámbulo ---
% \include da problemas antes del begin document. \input es seguro.
\input{Latex/Macros/MacroFile1}

% --- Definiciones de seguridad (por si la clase falla al cargar alguna variable) ---
\providecommand{\facultad}[1]{\def\lafacultad{#1}}
\providecommand{\escudofacultad}[1]{\def\elescudo{#1}}
\providecommand{\degree}[1]{\def\elgrado{#1}}
\providecommand{\director}[1]{\def\eldirector{#1}}
\providecommand{\lugar}[1]{\def\ellugar{#1}}
\providecommand{\degreedate}[1]{\def\lafecha{#1}}

% --- Configuración para bloques de código de R (Pandoc) ---

%%%%%%%%%%%%%%%%%%%%%%%%%%%%%%%%%%%%%%%%%%%%%%%%%%%%%%%%%%%%%%%%%%%%%%%%%%%%%%%%
%                         DATOS (Conectados a RMarkdown)                       %
%%%%%%%%%%%%%%%%%%%%%%%%%%%%%%%%%%%%%%%%%%%%%%%%%%%%%%%%%%%%%%%%%%%%%%%%%%%%%%%%
\title{Análisis Descriptivo de la Productividad Agrícola en Distintas
Regiones de la India Durante el Periodo 1997-2009}
\author{Asly Caputo \\ Aymara Andersen}

% Lógica para Facultad: Si la defines en el YAML la usa, si no, usa la default

\facultad{Facultad de Ciencias Económicas y Sociales \\ Escuela de
Estadistica y Ciencias Actuariales}
  \escudofacultad{Latex/Classes/Escudos/faces}

\degree{Computación I}
\director{Jesus Ochoa \\ Oliver Triveño}
\degreedate{Marzo 2026} 
\lugar{Caracas}

\portadatrue 

% Metadatos PDF
\keywords{}
\subject{}

% --- CORRECCIÓN PARA PANDOC NUEVO (Imágenes) ---
\makeatletter
\providecommand{\pandocbounded}[1]{#1}
\makeatother


%%%%%%%%%%%%%%%%%%%%%%%%%%%%%%%%%%%%%%%%%%%%%%%%%%%%%
%                   DOCUMENTO                       %
%%%%%%%%%%%%%%%%%%%%%%%%%%%%%%%%%%%%%%%%%%%%%%%%%%%%%
\begin{document}

\maketitle

%%%%%%%%%%%%%%%%%%%%%%%%%%%%%%%%%%%%%%%%%%%%%%%%%%%%%
%                  PRÓLOGO                          %
%%%%%%%%%%%%%%%%%%%%%%%%%%%%%%%%%%%%%%%%%%%%%%%%%%%%%
\frontmatter

% Puedes agregar dedicatorias desde el YAML si quieres, o dejarlas fijas aquí

%%%%%%%%%%%%%%%%%%%%%%%%%%%%%%%%%%%%%%%%%%%%%%%%%%%%%
%                   ÍNDICES                         %
%%%%%%%%%%%%%%%%%%%%%%%%%%%%%%%%%%%%%%%%%%%%%%%%%%%%%
\setcounter{secnumdepth}{3}
\setcounter{tocdepth}{3}

\tableofcontents



\cleardoublepage

  \cleardoublepage       % Empieza en página derecha (estándar en tesis)
  \chapter*{Resumen}     % Crea el título "Resumen" sin numeración (Capítulo *)
  \addcontentsline{toc}{chapter}{Resumen} % (Opcional) Lo añade al índice
  
  En el presente estudio se realizará un análisis estadístico
descriptivo del sector agrícola en la India, una de las economías a
nivel de producción de cultivo más relevantes a nivel global. El
objetivo principal es evaluar el comportamiento de la producción, el
rendimiento agrícola y el área cultivada en los diversos estados de la
India en el periodo comprendido entre 1997-2009 Utilizando el conjunto
de datos ``Agricultural Crop Yield in Indian States'' y herramientas de
programación en R y Python se realizará investigación que nos permitirá
decir la variabilidad del sector agrícola teniendo en cuenta factores
geográficos y los tipos de cultivos El estudio analiza una muestra de
2223 datos, enfocándose en variables como la Producción (Production),
área (Area), y el rendimiento (Yield) además de factores explicativos
Por último, en los resultados se busca identificar las regiones con
mayor productividad, la evolución histórica de los rendimientos y la
relación entre el tipo de cultivo y la extensión de tierra
utilizada.             % Aquí se pega el texto que escribiste en el YAML
  
  \cleardoublepage

%%%%%%%%%%%%%%%%%%%%%%%%%%%%%%%%%%%%%%%%%%%%%%%%%%%%%
%                   CONTENIDO (El Corazón)          %
%%%%%%%%%%%%%%%%%%%%%%%%%%%%%%%%%%%%%%%%%%%%%%%%%%%%%
\mainmatter
\def\baselinestretch{1.5}

% --- CAMBIO CRÍTICO 2: Aquí RMarkdown inyecta todo tu texto ---
\chapter*{Introducción}

La agricultura es una de las actividades económicas más cruciales en la
economía de cualquier país produciendo alimentos, materias primas y
generando empleos a gran parte de la población. Asimismo, la India
alberga una de las economías agrícolas más grandes del mundo con el
7,5\% de la producción agrícola a nivel global, convirtiéndose en uno de
los principales países productores agrícolas del mundo.\\

A causa de esto, resulta de interés investigar estadísticamente el
comportamiento del sector agrícola relacionado con la producción y
rendimiento de varios cultivos en la India para responder algunas
interrogantes de interés como:\\

\textbf{1.} ¿Cuáles son las regiones con mayor y menor producción
agrícola?\\

\textbf{2.} ¿Cómo ha evolucionado la producción y rendimiento a través
del tiempo?\\

\textbf{3.} ¿Cómo afecta el tipo de cultivo en el área cultivada?\\

Por eso, el presente estudio tiene como finalidad realizar un análisis
estadístico descriptivo del conjunto de datos ``rendimiento de los
cultivos agrícolas en los estados de la India'' que abarca datos
agrícolas de múltiples cultivos plantados en varios estados de la India
desde 1997 hasta 2020. Haciendo uso de herramientas para el análisis de
datos del lenguaje de programación R y de técnicas de visualización de
datos.\\

\section*{Estructura de Informe}

\begin{itemize}

    \item \textbf{Capítulo I}, se define el enfoque y propósito del estudio, planteando la problemática del estudio, formulando las preguntas de la investigación, determinando el contexto, límites y objetivos del problema a investigar.
    
    \item \textbf{Capítulo II}, se establecen las bases teóricas necesarias para entender cómo abordar el problema de estudio, definiendo conceptos primordiales para la investigación y las técnicas estadísticas necesarias para la resolución del problema.
    
    \item \textbf{Capítulo III}, se describe la metodología utilizada para el análisis de los datos del estudio, las técnicas usadas para la validación de los datos.
    
    \item \textbf{Capítulo IV}, se realiza la exposición de los resultados del análisis, respondiendo a las interrogantes de la investigación planteadas en el capítulo I haciendo uso de distintas técnicas de visualización  de datos como tablas y gráficos.
    
    \item \textbf{Capítulo V}, se muestran las conclusiones y recomendaciones basadas en el análisis de los resultados expuestos.
    
\end{itemize}

Este estudio busca entender las características del sector agrícola de
uno de los principales países productores agrícolas del mundo,
analizando su avance y evolución a través del tiempo.

\chapter{Planteamiento del Problema}

La agricultura ha sido un sector económico de la India que ha ido
evolucionando a través del tiempo, siendo actualmente uno de los países
con mayor producción agrícola a nivel global, lo que lo hace uno de los
sectores económicos más relevantes del país. Sin embargo, la agricultura
también es un sector con una alta variabilidad cuando hablamos de
rendimiento, dependiendo de factores como el tipo de cultivo sembrado y
la zona geográfica en la que se encuentra, por lo que obtener una alta
producción requiere de un análisis exhaustivo para obtener los mejores
resultados.\\

A nivel económico, la variabilidad del sector agrícola puede suponer un
riesgo significativo a la hora de invertir en este sector, por lo que es
necesario un estudio como mínimo la compatibilidad entre la región y el
tipo de cultivo escogidos, así como su evolución a través del tiempo.\\

Por eso, resulta de interés analizar estadísticamente el comportamiento
y rendimiento de la producción agrícola en la India a nivel general,
temporal y geográfico, ya que, independientemente de su gran nivel de
producción, es necesario saber la distribución del sector para definir
concretamente cual es el sitio más óptimo de inversión.\\

A partir de esto, se plantean las siguientes preguntas de
investigación:\\

\textbf{1.} ¿Cuáles son las regiones con mayor y menor producción
agrícola?\\

\textbf{2.} ¿Cómo ha evolucionado la producción y rendimiento a través
del tiempo?\\

\textbf{3.} ¿Cómo afecta el tipo de cultivo en el área cultivada?\\

\section{Justificación}

El presente estudio tiene como justificación la necesidad de conocer e
identificar estadísticamente la rentabilidad agrícola de una determinada
zona únicamente con los factores básicos que afectan la producción
agrícola, sin necesidad de tener un conocimiento profundo del área, lo
que permite la toma rápida de mejores decisiones.\\

Desde un punto estadístico, el explorar y analizar la variabilidad de
este sector ayuda a manejar técnicas del análisis descriptivo, haciendo
uso de herramientas tecnológicas como lo son los lenguajes R y Python.

\section{Objetivos}

\subsection{Objetivo General}

Realizar un análisis descriptivo de la producción, área cultivada y
rendimiento agrícola de las regiones de la India durante el periodo
1997-2009.

\subsection{Objetivos Específicos}

\begin{itemize}
  \item{Comparar la producción, área y rendimiento agrícolas total de la India durante todo el periodo de estudio.}
  \item{ Identificar los cultivos con mayor área cultivada.}
  \item{Identificar el impacto que tiene el área de cultivo según su tipo.}
  \item{Evaluar los cambios en la evolución histórica de la producción, área y rendimiento agrícola durante el periodo 1997-2009.}
  \item{Cuantificar la relación de la producción de las regiones de la India y los tipos de cultivos.}
\end{itemize}

\section{Cobertura}

\subsection{Horizontal}

La cobertura horizontal abarca las características de cada cultivo
registrado y distribuidas en 7 columnas, donde se consideran las
dimensiones descriptivas del cultivo y las métricas del rendimiento
junto con su relación entre diversos factores agronómicos.\\

\begin{table}[ht] % 
\centering
\caption{Variables consideradas en el estudio} 
\resizebox{12cm}{!} {
\begin{tabular}{cll}
  \hline
  \textbf{TIPO} & \textbf{VARIABLE} & \textbf{DESCRIPCIÓN} \\ 
  \hline
  Categórica & Crop & Nombre del cultivo. \\ 
  Númerica & Crop\_Year & El año en el que se sembró el cultivo. \\ 
  Categórica & State & Estado de la India donde se sembró el cultivo. \\ 
  Númerica & Area & Superficie total de tierra (en hectáreas) cultivada. \\ 
  Númerica & Production & Cantidad de producción agrícola (en toneladas métricas) \\ 
  Categórica & Crop\_Type & Tipo de cultivo sembrado. \\ 
  Númerica & Yield & Rendimiento calculado del cultivo (producción por unidad de área). \\ 
   \hline
\end{tabular}
}
\end{table}

\subsection{Vertical}

La cobertura vertical del trabajo de investigación abarca un total de
2237 observaciones (registros) donde cada fila representa un cultivo
específico y sus características correspondientes, con un alcance
temporal que abarca desde el año 1997 hasta el año 2009.

\section{Periodo de Referencia}

La fuente de datos a estudiar corresponde al repositorio abierto de
Kaggle (``Agricultural Crop Yield in Indian States Dataset''). Donde se
toma como período de referencia los registros pertenecientes al
intervalo de tiempo 1997-2009.

\section{Variables}

Para este estudio, se tomaron en cuenta como variables principales
Production, Area y Yield, las cuales representan la toneladas
producidas, las hectáreas cultivadas y su producción por área
respectivamente.\\

También se tomaron como variables explicativas:

\begin{enumerate}
  \item Variables Temporales: \textbf{Crop\_Year}
  \item Variables Categóricas: \textbf{Crop}
  \item Variables Geográficas: \textbf{State}
\end{enumerate}

\chapter{Marco Teórico}

\section{Bases Teóricas}

En este capítulo se desarrollarán los conceptos y fundamentos
estadísticos necesarios para realizar el análisis estadístico necesario
para cubrir los objetivos planteados anteriormente. Se abordarán
conceptos como la tipología de cultivos, regiones agrícolas de la India
y las medidas fundamentales para la Estadística descriptiva Univariante
y Bivariante.

Estos conceptos son necesarios para comprender la metodología y
resultados de análisis de este estudio y asociarlos al entorno agrícola
de la investigación.\\

\subsection{Tipología de Cultivos}

La India diversifica su producción para satisfacer tanto el mercado
interno como el global, en este estudio se analizan 50 variedades de
cultivos, por lo que agruparlos puede resultar útil para algunos
análisis específicos:

\begin{itemize}
 \item{\textbf{Cultivos Alimentarios:} Cereales y legumbres que forman la dieta básica. La India es líder mundial en la producción de legumbres (fuente principal de proteína para su población mayoritariamente vegetariana).}
 \item{\textbf{Cultivos de Plantación:} Introducidos en gran medida durante la época colonial, se cultivan en grandes fincas para exportación. Destacan el té (Assam y Darjeeling), el café (Karnataka) y el caucho (Kerala).}
 \item{\textbf{Cultivos Industriales o Comerciales:} Materias primas para la industria manufacturera. La caña de azúcar para el azúcar y el etanol, y el algodón para la potente industria textil india.}
 \item{\textbf{Horticultura:} Incluye frutas, flores y especias. La India es el mayor productor mundial de especias (pimienta, cardamomo, cúrcuma) y mangos.}
\end{itemize}

\hfill\break

\subsection{Regiones Agrícolas de la India}

El país se puede dividir en zonas productivas clave, para el estudio,
saber las zonas productivas nos da un indicio de cuáles cultivos pueden
tener máyor producción en algunos estados que otras:

\begin{itemize}

    \item{\textbf{Región del Noroeste (Punjab, Haryana):} Es la zona más tecnificada y rica. Produce la mayor parte del trigo y arroz del país.}
    \item{\textbf{Valle del Ganges y Brahmaputra:} Suelos aluviales extremadamente fértiles. Es el reino del arroz, el yute y la caña de azúcar.}
    \item{\textbf{Cinturón de Suelo Negro (Maharashtra, Gujarat):} Conocido como la "región del algodón". El suelo retiene muy bien la humedad, ideal para cultivos de fibra.}
    \item{\textbf{Zonas de Plantación del Sur y Noreste:} Las colinas de los Ghats Occidentales y Assam, donde el clima húmedo y la altitud permiten el cultivo de té, café y especias}
    \item{\textbf{Zonas Áridas (Rajasthan):} Se especializan en mijos y semillas oleaginosas que requieren mínima agua.}

\end{itemize}

\hfill\break

\subsection{Estadística Descriptiva Univariante}

La estadística descriptiva univariante se centra en el análisis de
datos, haciendo uso de herramientas para describir, tabular, representar
y visualizar el comportamiento de una variable.

Para este estudio, será necesario entender el comportamiento individual
de la producción, el área cultivada y el rendimiento anual utilizando
técnicas de la estadística descriptiva univariante.\\

\subsection*{Medidas de Posición}

Las medidas de posición, también llamadas fractiles o cuantiles, son
utilizadas para facilitar la localización de un dato respecto de un
conjunto de ellos. Para este estudio, se utilizarán:

\begin{itemize}
\item{\textbf{Cuartiles:} Valores que dividen la serie en 4 partes iguales, cada una de ellas con un 25\% de los datos.}

$$P_k = \frac{k(n + 1)}{4}$$

\item{\textbf{Mediana:} Promedio no matemático que divide a la serie de datos de un estudio, en dos partes iguales, que representan cada una el 50\% de los datos.}

$$Me = \left(\frac{n+1}{2}\right)$$
\end{itemize}

\hfill\break

\subsection*{Medidas de Tendencia Central}

Son casos especiales de las medidas de posición, debido a su ubicación
en la zona central de la distribución. Estas medidas, también son
denominadas promedios.

\begin{itemize}
\item{\textbf{Media Aritmética:} Es el valor resultante más “representativo” de la serie de datos, el punto de equilibrio, el centro de gravedad de la serie.}

$$\bar{x} = \frac{\sum_{i=1}^{n} x_i}{n}$$
\item{\textbf{Mediana:} Promedio no matemático que divide a la serie de datos de un estudio, en dos partes iguales, que representan cada una el 50\% de los datos.}

 $$Me = \left(\frac{n+1}{2}\right)$$

\item{\textbf{Moda:} Es un promedio no matemático representado por el valor con mayor número de ocurrencias dentro de la serie de datos.}
\end{itemize}

\hfill\break

\subsection*{Medidas de Dispersión}

Tienen por objeto medir la magnitud de los desvíos de los valores de la
variable con respecto al valor central de la distribución, es decir,
definen cuán semejante o cuán distinto son cada uno de los valores, de
la variable con respecto al valor central.

\begin{itemize}
\item{\textbf{Rango Intercuartílico:} constituye la máxima diferencia posible entre dos valores entre el 50\% de los datos centrales.}

$$IQR = Q_3 - Q_1$$

\item{\textbf{Varianza:} Es el promedio cuadrático de los desvíos de los datos respecto de la media aritmética.}

$$\sigma^2 = \frac{\sum_{i=1}^{N} (x_i - \mu)^2}{N}$$

\item{\textbf{Desviación Típica:} Es la raíz cuadrada positiva de la varianza. Constituye el índice de dispersión más usado y más confiable.}

$$s = \sqrt{s^2} = \sqrt{\frac{\sum_{i=1}^{n} (x_i - \bar{x})^2}{n}}$$

\item{\textbf{Coeficiente de Variación de Pearson:} Es el cociente entre la desviación estándar y la media aritmética, multiplicado por cien. Es el coeficiente de dispersión relativa de más amplio uso.}

$$CV =  \frac{s}{\bar{x}} \times 100\%$$
\end{itemize}

\hfill\break

\subsection{Estadística Descriptiva Bivariante}

La estadística descriptiva bivariante analiza la relación entre dos
variables donde se busca resumir y describir la información conjunta de
ambos conjuntos de datos, para identificar patrones y tendencias

En este estudio resulta de interés conocer a fondo el tipo de relación
que hay entre cada una de las variables, tanto numéricas como
categóricas, para saber la influencia que tienen una sobre a otra.\\

\subsection*{Medidas de Dispersión Conjunta}

\begin{itemize}
\item{\textbf{Coeficiente de Correlación de Pearson:} Es una medida que da cuenta del sentido, dirección e intensidad o magnitud de la relación lineal entre dos variables cuantitativas.}

$$r = \frac{\sum_{i=1}^{n} (x_i - \bar{x})(y_i - \bar{y})}{n\cdot\sqrt{\sum_{i=1}^{n}) (x_i - \bar{x})^2} \cdot \sqrt{\sum_{i=1}^{n} (y_i - \bar{y})^2}}$$
\end{itemize}

\hfill\break

\subsubsection*{Método de Detección de Datos Atípicos}
\subsubsection*{Prueba de Tukey}

verifica múltiples valores atípicos en ambos lados, categorizados como
valores ``externos'' o ``muy alejados''. En este caso los límites son:

\begin{itemize}
\item{\textbf{Límites internos:} Se consideran atípicas leves aquellas observaciones alejadas del 1er y 3er cuartil, respectivamente, más 1,5 veces el Recorrido Intercuartílico.}

$$X_i \le Q_1 - 1,5RI \quad \text{o} \quad X_i \ge Q_3 + 1,5RI$$

\item{\textbf{Límites externos:} Se consideran atípicas extremas aquellas observaciones alejadas del 1er y 3er cuartil, respectivamente, más de 3 veces el Recorrido Intercuartílico.}

$$X_i \le Q_1 - 3RI \quad \text{o} \quad X_i \ge Q_3 + 3RI$$
\end{itemize}

\hfill\break

\subsection{Series de Tiempo}

Una serie de tiempo se conoce como una representación visual de datos
donde se observa la evolución, bien sea de días, meses o años. A
diferencia de otros conjuntos de datos, en las series de tiempo es
sumamente importante que vayan en orden cronológico ya que los valores
menores suelen influir en los valores futuros

En la agricultura se puede utilizar la serie de tiempo para medir el año
de cosecha y su evolución

\chapter{Marco Metódico}

En este capítulo se describe el proceso de selección de aquellos
registros necesarios para cumplir con los objetivos planteados en el
capítulo I. Para esto se hará uso de herramientas para el análisis de
datos como R y Pandas (librería de Python).\\

\section{Bases metódicas utilizadas en el trabajo de investigación}

En la presente investigación fijamos una variable cuantitativa a nivel
descriptivo, ya que, se busca caracterizar la situación agrícola
procesando datos numéricos de la producción, áreas cultivadas y
rendimiento agrícola en la India. Dicho esto se utilizó bases metódicas
estadísticas integrando datos cuantitativos y cualitativos.\\

\subsubsection{Verificación de Datos Nulos}

Con la ayuda de Tidyverse utilizando librerías como los son (``dplyr'' y
``tidyr''), se escanea el dataset en busca de valores nulos que puedan
afectar el análisis. Dando como resultado la ausencia de datos nulos en
todas las columnas del dataset original.

Por lo tanto, podemos decir que no hay datos nulos en el dataset a
trabajar.

\begin{table}[!h]
\centering
\caption{\label{tab:unnamed-chunk-1}Valores Nulos por Columna}
\centering
\fontsize{10}{12}\selectfont
\begin{tabular}[t]{rrrrrr}
\toprule
Crop\_Year & State & Crop & Production & Area & Yield\\
\midrule
\cellcolor{gray!10}{0} & \cellcolor{gray!10}{0} & \cellcolor{gray!10}{0} & \cellcolor{gray!10}{0} & \cellcolor{gray!10}{0} & \cellcolor{gray!10}{0}\\
\bottomrule
\end{tabular}
\end{table}

\subsubsection{Selección de Datos Necesarios}

Como se mencionó en el Capítulo I, se tomarán los registros anuales
entre el periodo de 1997 hasta el 2009. Esto implica que las variables
``Crop\_Year'' y ``Season'' serán filtradas según esas condiciones.\\

Aparte, se seleccionan únicamente las columnas que contengan las
variables necesarias para el estudio, es decir:\\

Crop, Crop\_Year, Production, Area y Yield.\\

\subsubsection{Categorización de Cultivos}

Para facilitar el análisis, se agruparon sectores de cultivos creando
nueva información colocada en etiquetas. Un cultivo como ``Arroz'' y
``Frijoles'' se colocó en la nueva Variable (Crop\_Year) como ``Cereales
y Granos'' , permitiendo comparar el crecimiento de cereales básicos con
legumbres o cultivos comerciales.\\

\subsubsection{Control de Outliers}

Para visualizar correctamente la distribución de alguna variable
númerica como Production, Area ó Yield, es necesario identificar la
cantidad de valores atípicos de cada una, para ello se utilizó el método
del rango intercuartílico, con el fin de saber la cantidad de outliers
de cada una y el porcentaje que ocupan en la variable.\\

\begin{table}[!h]
\centering
\caption{\label{tab:unnamed-chunk-4}Cantidad de Datos Atípicos por Columna}
\centering
\fontsize{10}{12}\selectfont
\begin{tabular}[t]{lrrr}
\toprule
Variable & Registros & Total\_Atipicos & Porcentaje\\
\midrule
\cellcolor{gray!10}{Produccion} & \cellcolor{gray!10}{2238} & \cellcolor{gray!10}{435} & \cellcolor{gray!10}{19.44}\\
Area & 2238 & 286 & 12.78\\
\cellcolor{gray!10}{Rendimiento} & \cellcolor{gray!10}{2238} & \cellcolor{gray!10}{337} & \cellcolor{gray!10}{15.06}\\
\bottomrule
\end{tabular}
\end{table}

\chapter{Análisis de Resultados}

\{\small Este capítulo presenta los resultados obtenidos al aplicar las
metodologías descritas en el capítulo anterior, así como del análisis
estadístico descriptivo de las variables estudiadas y las ventajas y
desventajas al aplicar dichos tests.\}

\subsubsection*{MEDIDAS DE TENDENCIA CENTRAL Y DISPERCIÓN}

\begin{table}[!h]
\centering
\caption{\label{tab:unnamed-chunk-5}Medidas de Tendencia Central y Dispersión}
\centering
\fontsize{10}{12}\selectfont
\begin{tabular}[t]{lrrr}
\toprule
Indicador & Media & Mediana & Desviación\_Tipica\\
\midrule
\cellcolor{gray!10}{Production} & \cellcolor{gray!10}{68741614.34} & \cellcolor{gray!10}{24762.00} & \cellcolor{gray!10}{537174070.52}\\
Area & 67723.38 & 7476.50 & 238939.79\\
\cellcolor{gray!10}{Yield} & \cellcolor{gray!10}{336.74} & \cellcolor{gray!10}{3.84} & \cellcolor{gray!10}{1732.53}\\
\bottomrule
\end{tabular}
\end{table}

Como podemos observar existe una gran diferencia entre los promedios y
los valores centrales. En las tres variables (Producción, Área y
Rendimiento), la Media es muchísimo más alta que la Mediana. Esto
significa que tienes unos pocos casos con números gigantescos que
``inflan'' el promedio, aunque la mayoría de los datos son valores
pequeños. Por ejemplo, en la Producción, el promedio es de 68 millones,
pero el valor central real (mediana) es de apenas 24 mil. La desviación
Típica es extremadamente alta, lo que nos dice que los datos están muy
dispersos y no se parecen entre sí; hay mucha desigualdad en la
eficiencia y el tamaño de las tierras.

\subsubsection*{REGIONES CON MAYOR Y MENOR PRODUCCIÓN}

\begin{center}\includegraphics[width=0.9\linewidth,]{informe_files/figure-latex/unnamed-chunk-6-1} \end{center}

Al observar el grafico, tenemos como resultado una distribución
desigual. Las regiones del norte y el centro del país presentan una
actividad muy escasa, mientras que el sur y zonas montañosas del país
aparecen con color más intenso, lo que nos indica que la producción
agrícola es más alta.

Respondiendo así a la pregunta principal ¿Cuáles son las regiones con
mayor y menor producción agrícola?.

Los estados con mayor producción se encuentran al sur de la India,
teniendo a Kerala como el estado con más de 74000 millones; mientras que
Jammu and Kashmir se sitúa como el estado con menor producción agrícola
con un total de 63330.\\

\subsubsection*{¿CÓMO HA EVOLUCIONADO LA PRODUCCIÓN Y EL RENDIMIENTO A TRAVÉS DEL TIEMPO?}

\begin{center}\includegraphics[width=0.9\linewidth,]{informe_files/figure-latex/unnamed-chunk-7-1} \end{center}

El gráfico de evolución temporal, el cual indica la respuesta a nuestro
objetivo numero 2, nos muestra un crecimiento constante. Desde 1997
hasta el año 2009, no ha decaído en ningún momento. Aunque se muestran
pequeños saltos entre años (subidas y bajadas) probablemente causados
por el clima, la tendencia general es muy positiva; demostrando que la
India no solo ha mantenido el ritmo de producción, sino que produce cada
vez más.\\

\subsubsection*{¿CÓMO AFECTA EL TIPO DE CULTIVO EN EL ÁREA CULTIVADA?}

\begin{center}\includegraphics[width=0.9\linewidth,]{informe_files/figure-latex/unnamed-chunk-8-1} \end{center}

En el análisis de las barras observamos que la papa, el banano y la caña
de azúcar son los cultivos que lideran el volumen de producción general,
afectando de manera positiva en cuanto a su rendimiento. Destacan por su
eficiencia por área, ya que producen mucho en poco espacio. Al enfocarse
en estos cultivos esto le permite al país tener más ganancias y asegurar
alimento para la población.\\

\subsubsection*{MATRIZ DE CORRELACIÓN DEL ÁREA, PRODUCCIÓN, Y RENDIMIENTO AGRÍCOLA}
\begin{table}[!h]
\centering
\caption{\label{tab:unnamed-chunk-9}Matriz de Correlación de Producción, Rendimiento y Área}
\centering
\fontsize{10}{12}\selectfont
\begin{tabular}[t]{lrrr}
\toprule
  & Production & Area & Yield\\
\midrule
\cellcolor{gray!10}{Production} & \cellcolor{gray!10}{1.00} & \cellcolor{gray!10}{0.28} & \cellcolor{gray!10}{0.54}\\
Area & 0.28 & 1.00 & 0.10\\
\cellcolor{gray!10}{Yield} & \cellcolor{gray!10}{0.54} & \cellcolor{gray!10}{0.10} & \cellcolor{gray!10}{1.00}\\
\bottomrule
\end{tabular}
\end{table}

Relacion entre produccion y rendimiento (r = 0.543): La relación más
clara está entre la producción y el rendimiento. Esto nos dice que el
volumen total de la cosecha depende principalmente de qué tan bien se
aprovecha cada hectárea. Básicamente, si quieres producir más, la clave
es mejorar la eficiencia de tus cultivos.

Relacion entre produccion y area (r = 0.280): La relación entre la
Producción y el Área es baja. Esto significa que sembrar más tierra no
provoca automáticamente un gran salto en la producción total. En la
práctica, es mucho más efectivo trabajar con lo que ya tienes que
simplemente tratar de expandirse a más hectáreas.

Relación entre terreno y rendimiento (r = 0.102): Casi no hay relación
entre el tamaño del terreno y el rendimiento. Esto indica que la
eficiencia es la misma sin importar si la parcela es pequeña o grande.
No estás ganando beneficios extra por ser más grande, pero tampoco estás
perdiendo calidad al crecer.

\chapter{Conclusiones y Recomendaciones}

Después de realizar el análisis de datos en el presente estudio y
observar la evolución de la producción podemos concluir lo siguiente:

\textbf{1.-} En el periodo del año 1997 hasta el 2009 la agricultura
tuvo una evolución constante, aunque tuvo altas y bajas a causa
principalmente del clima al pasar los años se visualizó que nunca dejó
de crecer su producción.~

\textbf{Recomendación} Tener en cuenta la temporada en la que se cultiva
para que el clima no sea un factor en contra y así poder tener mejor
rendimiento con respecto a la producción.\\

\textbf{2.-} Se determinó que los sectores con mayor producción son
aquellos ubicados al sur de la India, siendo esos los estados más
montañosos. Esto nos permite concluir que la producción está ligada a la
geografía, concentrándose en las zonas donde hay mejor tierra, su clima
es más húmedo y cuenta con zonas más altas, lo que permite una mejor
cosecha.\\

\textbf{Recomendación} Priorizar los cultivos que se producen en esta
área. Sin embargo, tomando en cuenta que son sitios húmedos y suelen
contar con la existencia de hongos, se recomienda tomar las medidas
necesarias para que los cultivos no sufran ningún tipo de daño.\\

\textbf{3.-} Por último se observó que los cultivos con mayor producción
(papa, banana, cebolla) son los cultivos que tienen mayor impacto en el
área de alimentos y economía.\\

\textbf{Recomendación} Enfocarse principalmente en estos cultivos a la
hora de su producción, de manera que logre ser considerado una fuente
principal de ingresos y alimentos.\\

\appendix
\renewcommand{\thechapter}{\Alph{chapter}}

\chapter{Anexo A}

\begin{center}\includegraphics[width=0.9\linewidth,]{informe_files/figure-latex/unnamed-chunk-10-1} \end{center}

\chapter{Anexo B}

\begin{center}\includegraphics[width=0.9\linewidth,]{informe_files/figure-latex/unnamed-chunk-11-1} \end{center}

\chapter{Anexo C}

\begin{center}\includegraphics[width=0.9\linewidth,]{informe_files/figure-latex/unnamed-chunk-12-1} \end{center}

\chapter{Anexo D}

\begin{center}\includegraphics[width=0.9\linewidth,]{informe_files/figure-latex/unnamed-chunk-13-1} \end{center}

\chapter{Anexo E}

\begin{center}\includegraphics[width=0.9\linewidth,]{informe_files/figure-latex/unnamed-chunk-14-1} \end{center}

\chapter{Bibliografia}

\textbf{1.} Pinto, S. \& Salomón, M. (2015): Análisis descriptivo
univariante, material docente. Caracas: Escuela de Estadística y
Ciencias Actuariales - UCV.

\textbf{2.} Jairo A. Ayala Godoy RPubs by RStudio

\textbf{3.} Tipos de agricultura en India GeeksforGeek
\url{https://www.geeksforgeeks.org} › types-of-agriculture

%%%%%%%%%%%%%%%%%%%%%%%%%%%%%%%%%%%%%%%%%%%%%%%%%%%%%
%                   REFERENCIAS                     %
%%%%%%%%%%%%%%%%%%%%%%%%%%%%%%%%%%%%%%%%%%%%%%%%%%%%%
% Mantenemos la estructura natbib de tu plantilla
\renewcommand{\bibname}{Lista de Referencias}
\bibliographystyle{apalike} 
\bibliography{bibliografia.bib} 

%%%%%%%%%%%%%%%%%%%%%%%%%%%%%%%%%%%%%%%%%%%%%%%%%%%%%
%                   APÉNDICES                       %
%%%%%%%%%%%%%%%%%%%%%%%%%%%%%%%%%%%%%%%%%%%%%%%%%%%%%

\end{document}
